% Generated draft: requires mathematical and source review.
\documentclass[11pt]{article}
\usepackage[margin=25mm]{geometry}
\usepackage{fontspec}
\setmainfont{Latin Modern Roman}
\usepackage{amsmath,amssymb,mathtools,unicode-math}
\setmathfont{Latin Modern Math}
\usepackage{enumitem,needspace,xurl,hyperref}
\setlength{\parindent}{0pt}
\setlength{\parskip}{5pt}
\setlength{\emergencystretch}{3em}
\newcommand{\N}{\mathbb{N}}
\newcommand{\Z}{\mathbb{Z}}
\newcommand{\Q}{\mathbb{Q}}
\newcommand{\R}{\mathbb{R}}
\newcommand{\C}{\mathbb{C}}
\newcommand{\F}{\mathbb{F}}
\newcommand{\A}{\mathbb{A}}
\newcommand{\PP}{\mathbb{P}}
\newcommand{\E}{\mathbb{E}}
\newcommand{\eps}{\varepsilon}
\newcommand{\dd}{\,\mathrm{d}}
\begin{document}
\section*{UnsolvedMath ID: 701710}
\textbf{Source problem code:} catalogue entry
\subsection{Definitions and mathematical statement}
For a finite group \(G\), let \(|G|\) denote its order and let
\[
\operatorname{Cl}(G)=\bigl\{\{xgx^{-1}:x\in G\}:g\in G\bigr\},\qquad k(G)=|\operatorname{Cl}(G)|
\]
be its set and number of conjugacy classes. The supplied source attributes to Burnside the result that \(16\mid(|G|-k(G))\) for every finite group \(G\) of odd order.

The question is whether
\[
\forall m\in\Z_{>0}\ \exists\text{ a finite nonabelian group }G(m)
\quad\text{such that}\quad
|G(m)|\text{ is odd},\qquad |G(m)|-k(G(m))=16m.
\]
Here nonabelian means that there exist \(a,b\in G(m)\) with \(ab\ne ba\).
\subsection{Short English statement}
Does every positive multiple of \(16\) occur as the order minus the number of conjugacy classes of a finite nonabelian group of odd order? This is the question proposed by D. MacHale in Kourovka Notebook Problem 20.52.
\subsection{Sources}
\begin{itemize}
\item Local UnsolvedMath registry, permanent ID \(701710\), exact source code \texttt{catalogue entry}, title \emph{Kourovka Notebook Problem 20.52}. The supplied record labels the problem open; this status is not independently verified.
\item \emph{The Kourovka Notebook (Unsolved Problems in Group Theory)}, 20th issue (2022), original problem 20.52, proposed by D. MacHale. The supplied background also identifies the 21st issue edition (2026), \texttt{arXiv:1401.0300v48}. The collection record names E. I. Khukhro and V. D. Mazurov as editors.
\item The supplied statement attributes the divisibility result to Burnside but gives no separate bibliographic citation for that result. These references and attributions are supplied source information and have not been independently checked.
\end{itemize}
\end{document}
